% Generated draft: requires mathematical and source review.
\documentclass[11pt]{article}
\usepackage[margin=25mm]{geometry}
\usepackage{fontspec}
\setmainfont{Latin Modern Roman}
\usepackage{amsmath,amssymb,mathtools,unicode-math}
\setmathfont{Latin Modern Math}
\usepackage{enumitem,needspace,xurl,hyperref}
\setlength{\parindent}{0pt}
\setlength{\parskip}{5pt}
\setlength{\emergencystretch}{3em}
\newcommand{\N}{\mathbb{N}}
\newcommand{\Z}{\mathbb{Z}}
\newcommand{\Q}{\mathbb{Q}}
\newcommand{\R}{\mathbb{R}}
\newcommand{\C}{\mathbb{C}}
\newcommand{\F}{\mathbb{F}}
\newcommand{\A}{\mathbb{A}}
\newcommand{\PP}{\mathbb{P}}
\newcommand{\E}{\mathbb{E}}
\newcommand{\eps}{\varepsilon}
\newcommand{\dd}{\,\mathrm{d}}
\begin{document}
\section*{UnsolvedMath ID: 20000539}
\textbf{Source problem code:} catalogue entry
\subsection{Definitions and mathematical statement}
For $0<\eps<1/2$, real $\lambda$ and initially smooth compactly supported $f:\R\to\C$, use the reconstructed symmetric principal value
\[T_{\eps,\lambda}f(x)=\operatorname{p.v.}\int_\R f(x-t)e^{i\lambda|t|^\eps}\frac{\dd t}{t},\qquad M_\eps f(x)=\sup_{\lambda\in\R}|T_{\eps,\lambda}f(x)|.\]
\textbf{Original Question 12 (Guo).} Is there a universal constant $C>0$, independent of $\eps$, such that
\[\|M_\eps f\|_{L^2(\R)}\le C\|f\|_{L^2(\R)}\quad\text{for every }0<\eps<1/2?\]
This is a uniform family of maximal-in-$\lambda$ bounds, not a single maximal operator also taking a supremum in $\eps$. The exponent and principal-value/absolute-value symbols are reconstructed from damaged extraction; the registry marks the statement \texttt{reconstructed\_unverified}.

\textbf{Attached unreviewed endpoint reduction.} For fixed real $\mu$, choose the allowed scaling $\lambda=\mu/\eps$ and remove the scalar phase. The aid claims pointwise convergence on $C_c^1(\R)$,
\[e^{-i\mu/\eps}T_{\eps,\mu/\eps}f(x)\longrightarrow S_\mu f(x):=\operatorname{p.v.}\int_\R f(x-t)|t|^{i\mu}\frac{\dd t}{t}.\]
Consequently the universal bound would imply, with the same constant,
\[\left\|\sup_{\mu\in\R}|S_\mu f|\right\|_2\le C\|f\|_2.\]
The aid obtains this through finite parameter maxima and Fatou's lemma, then the appropriate parameter-continuity/density argument; it does not exchange an unproved infinite supremum with a limit. The logarithmic model is necessary, not claimed sufficient for the original estimate.

With Fourier convention $\widehat f(\xi)=\int e^{-ix\xi}f(x)\dd x$, the reported exact multiplier for $\xi\ne0$ is
\[\widehat{S_\mu f}(\xi)=2\Gamma(i\mu)\sinh(\pi\mu/2)\operatorname{sgn}(\xi)|\xi|^{-i\mu}\widehat f(\xi).\]
At $\mu=0$ its removable value is $-i\pi\operatorname{sgn}\xi$, the unnormalized Hilbert-transform multiplier. The gamma expression is interpreted by continuation at that parameter.

For every fixed $0<a<b<\infty$ and $M<\infty$, define annular convolution kernels
\[K_{\eps,\mu}(t)=\mathbf1_{\{a\le|t|\le b\}}\frac{\exp\bigl(i\mu(|t|^\eps-1)/\eps\bigr)}t,\qquad K_{0,\mu}(t)=\mathbf1_{\{a\le|t|\le b\}}\frac{|t|^{i\mu}}t.\]
The aid asserts $\sup_{|\mu|\le M}\|K_{\eps,\mu}*-K_{0,\mu}*\|_{L^2\to L^2}=O_{a,b,M}(\eps)$. This is fixed-annulus, bounded-$\mu$ control, not full-kernel operator-norm convergence uniformly over all parameters. The aid reports the universal AIM estimate as open; all reduction and multiplier assertions require independent review.
\subsection{Short English statement}
Guo asks for an $L^2$ bound uniform as the power exponent tends to zero. The aid identifies a necessary logarithmic-modulation maximal problem by scaling the phase coefficient like $1/\eps$, gives its exact multiplier, and quantifies convergence only on fixed annuli with bounded logarithmic parameter. It does not prove the requested universal bound.
\subsection{Sources}
\begin{itemize}
\item Local UnsolvedMath registry, permanent ID 20000539, source code catalogue entry; \texttt{aim-analysis-notes.json}, notes[70], tag \texttt{question}.
\item AIM workshop \emph{Carleson theorems and multilinear operators}, Question 12 (Guo), equation (0.17): \url{https://aimath.org/pastworkshops/multilinopsproblems.pdf}. Supplied verification describes comparison with the PDF and Guo's operator definition, but gives no separate bibliographic locator. No PDF was consulted here.
\item Attached machine-generated research attempt 1: result \texttt{reduction}, problem status at run \texttt{open}, rigor label \texttt{proved}, novelty confidence low. It is not peer reviewed and requires independent verification before citation. Original literature field empty.
\end{itemize}
\end{document}
